\begin{scanpage}{0213}{195}
\[\left.\begin{aligned}
\mathcal E^{(3)}(z_1)&=-\mathcal E_2^{(2)}\psi_-(L)e^{ikL},\\
\mathcal E^{(3)}(z_2)&=-\mathcal E_1^{(2)}\psi_+(L)e^{ikL}.
\end{aligned}\right\}\tag{31.09}\]
In the directions toward the opposite end of the conductor, these waves are equivalent to the radiation of an infinite single-wire line excited by the external fields
\[E_z^e=\mathcal E_1^{(3)}\cdot\delta(z-z_1),\quad\mathcal E_1^{(3)}=\left.\mathcal E^{(3)}(z_1)\right|_{\vartheta=0},\tag{31.10}\]
\[E_z^e=\mathcal E_2^{(3)}\cdot\delta(z_2-z),\quad\mathcal E_2^{(3)}=\left.\mathcal E^{(3)}(z_2)\right|_{\vartheta=\pi}.\tag{31.11}\]
Consequently, the quaternary waves again may be investigated as edge waves radiated by the semi-infinite conductors $(-\infty,z_2)$ and $(z_1,\infty)$ with their excitation by the external fields (31.10) and (31.11). Using the reciprocity principle, we easily obtain
\[\left.\begin{aligned}
E_\vartheta^{(4)}(z_1)=H_\varphi^{(4)}(z_1)&=\frac{\mathcal E^{(4)}(z_1)}{2\sin\vartheta\ln\dfrac{2i}{\gamma ka\sin\vartheta}}\frac{e^{ikR}}R e^{-ikz_1\cos\vartheta},\\
E_\vartheta^{(4)}(z_2)=H_\varphi^{(4)}(z_2)&=\frac{\mathcal E^{(4)}(z_2)}{2\sin\vartheta\ln\dfrac{2i}{\gamma ka\sin\vartheta}}\frac{e^{ikR}}R e^{-ikz_2\cos\vartheta},
\end{aligned}\right\}\tag{31.12}\]
where
\[\left.\begin{aligned}
\mathcal E^{(4)}(z_1)&=-\mathcal E_2^{(3)}\psi_-(L)e^{ikL},\\
\mathcal E^{(4)}(z_2)&=-\mathcal E_1^{(3)}\psi_+(L)e^{ikL}.
\end{aligned}\right\}\tag{31.13}\]
In a completely similar way, the $n^{\rm th}$ order edge waves
\[\left.\begin{aligned}
E_\vartheta^{(n)}(z_1)=H_\varphi^{(n)}(z_1)&=\frac{\mathcal E^{(n)}(z_1)}{2\sin\vartheta\ln\dfrac{2i}{\gamma ka\sin\vartheta}}\frac{e^{ikR}}R e^{-ikz_1\cos\vartheta},\\
E_\vartheta^{(n)}(z_2)=H_\varphi^{(n)}(z_2)&=\frac{\mathcal E^{(n)}(z_2)}{2\sin\vartheta\ln\dfrac{2i}{\gamma ka\sin\vartheta}}\frac{e^{ikR}}R e^{-ikz_2\cos\vartheta},
\end{aligned}\right\}\tag{31.14}\]
are found. Here
\end{scanpage}
